% Part: incompleteness
% Chapter: theories-computability
% Section: intepretability

\documentclass[../../../include/open-logic-section]{subfiles}

\begin{document}

\olfileid{inc}{tcp}{itp}

\olsection{Teori kang ing jerone $\Th{Q}$ bisa diinterpretasikake ora bisa diputusake}

Kita bisa ngiyatake asil-asil iki luwih adoh. Sacara informal,
interpretasi basa $\Lang{L_1}$ ing basa liyane
$\Lang{L_2}$ kalebu nemtokake semesta, simbol relasi, lan
simbol fungsi saka $\Lang{L_1}$ nganggo !!{formula}s ing
$\Lang{L_2}$. Sanajan kita ora bakal ngrembug kanthi rinci, definisi
iki bisa digawe presisi.

\begin{thm}
  Upamana $\Th{T}$ teori ing basa kang bisa nginterpretasikake basa
  aritmetika, kanthi cara saengga $\Th{T}$ konsisten karo interpretasi
  $\Th{Q}$. Mula $\Th{T}$ ora bisa diputusake. Yen $\Th{T}$
  mbuktekake interpretasi aksioma-aksioma $\Th{Q}$, mula ora ana
  ekstensi konsisten saka $\Th{T}$ kang bisa diputusake.
\end{thm}

Buktine mung owah-owahan cilik saka bukti teorema sadurunge;
sawijining conto tandhing bisa digunakake kanggo misahake $\Th{Q}$ lan
$\Th{\bar Q}$. $\Th{ZFC}$, teori himpunan Zermelo--Fraenkel kanthi
aksioma pilihan, bisa dianggep minangka landhesan aksiomatik kang cukup
kuwat kanggo nindakake akeh matematika lumrah. Ing $\Th{ZFC}$,
bilangan asli bisa ditemtokake, lan liwat interpretasi iki,
aksioma-aksioma $\Th{Q}$ bener. Mula kita nduweni

\begin{cor}
Ora ana ekstensi $\Th{ZFC}$ kang bisa diputusake.
\end{cor}

\begin{cor}
Ora ana ekstensi $\Th{ZFC}$ kang jangkep, konsisten, lan
!!{axiomatizable} kanthi komputabel.
\end{cor}

Basa $\Th{ZFC}$ mung nduweni siji relasi biner,
$\in$. (Malah, kesamaan uga ora perlu.) Mula kita nduweni

\begin{cor}
Logika orde siji kanggo saben basa kang nduweni simbol relasi biner
ora bisa diputusake.
\end{cor}

Asil iki uga lumaku kanggo saben basa kang nduweni rong simbol fungsi
uner, amarga simbol-simbol iki bisa digunakake kanggo nyimulasi simbol
relasi biner. Asil-asil mau pas ing watese: jebul logika orde siji kanggo
basa kang mung nduweni simbol relasi \emph{uner} lan paling akeh siji
simbol fungsi \emph{uner} bisa diputusake.

Ana siji cathetan tambahan. Kita ngerti yen himpunan ukara ing
basa $\Obj 0$, $'$, $+$, $\times$, $<$ kang bener ing model standar
ora bisa diputusake. Nyatane, $<$ bisa ditemtokake nganggo simbol-simbol
liyane, banjur $+$ bisa ditemtokake nganggo $\times$ lan $'$. Mula
himpunan ukara kang bener ing basa $\Obj 0$, $'$, $\times$ ora
bisa diputusake. Kosok baline, Presburger wis nuduhake yen himpunan
ukara ing basa $\Obj 0$, $'$, $+$ kang bener ing basa aritmetika
bisa diputusake. Nanging prosedure ora praktis sacara komputasional.

\end{document}
